\section{Introduction}

Modern network services depend on implementations of protocol standards such as
TLS, MQTT, and QUIC. These standards specify how endpoints authenticate peers,
negotiate cryptographic parameters, reject malformed messages, preserve session
state, and interoperate across independent
implementations~\cite{rfc8446,rfc9000,rfc9001,rfc9002,mqtt311,mqtt5}. Even when
the standard is clear, an implementation can still diverge in subtle
field-handling, state-checking, or error-path behavior. Such implementation
deviations are not merely benign compatibility quirks: they may cause an
implementation to accept a forbidden field value, skip a required
state-dependent check, reject a legal message, or expose different behavior
under a build option that deployers rarely test. The consequences range from
interoperability failures to downgrade opportunities and denial-of-service
conditions.

Checking these deviations is difficult because the two sides of the problem are
written in radically different languages. Standards are natural-language
documents whose normative requirements are distributed across field tables,
message descriptions, state-machine prose, compatibility notes, and exception
clauses. RFC-style normative language provides useful cues~\cite{rfc2119,rfc8174}, but
those cues do not by themselves identify executable code obligations.
Implementations are engineering artifacts with macros, configuration
flags, generated code, helper functions, callbacks, and optimized error paths.
The most important evidence for a single rule may be split across a parser, a
state transition, a validation helper, and a caller-imposed precondition. A
checker must therefore answer a consistency question rather than a mere
localization question: under the conditions described by a normative rule, does
the implementation enforce the required behavior on the relevant code paths?

Existing approaches cover only parts of this question. Formal methods can prove
strong properties once a formal model exists, but constructing such a model from
an evolving RFC or industry standard remains expensive and requires expertise.
Protocol fuzzing and conformance tests exercise implementations, but they often
use test oracles that are written separately from individual normative
requirements, making it difficult to tell which field-level rule a test
validates or leaves unchecked.
Static analysis can reason over code structure, yet it typically has no access
to the semantics of the standard. Recent large language models provide a new
opportunity: their ability to jointly interpret natural-language requirements
and source code makes it possible to bridge standards and implementations
without requiring a complete hand-written formal model. However, this
opportunity is useful for conformance checking only when the model is given
precise obligations and the code evidence needed to evaluate them.

This paper is based on a simple observation: although protocol standards are
written as prose, many checkable requirements are ultimately anchored in fields.
A standard may state that a field must be present only under a condition, that a
length must match another field, that a value must be derived from a negotiated
parameter, or that a malformed combination must trigger an alert. These rules
can be represented as field-centric assertions of the form \ruleform: under
condition $C$, field $f$ must satisfy action or constraint $A$. This
representation is not a full formal semantics of the protocol. Instead, it is a
practical intermediate layer for organizing evidence, aligning specifications
to code, and making uncertainty explicit when the current evidence is not
enough.

We present \tool, a framework for checking whether protocol implementations
satisfy field-level requirements in natural-language standards. \tool first
constructs a field space from message structures, extensions, and field
definitions, then extracts field-centric rules from normative prose. It next
maps each rule to code by combining LLM-based semantic ranking, identifier
normalization, Tree-sitter and language-native AST evidence, and deterministic
filtering. Rather than asking a model for a one-shot verdict, \tool performs
iterative consistency reasoning: the model must report supporting evidence,
missing evidence, and one of three outcomes: consistent, inconsistent, or
insufficient evidence. When the outcome is suspicious or under-supported, a
test-validation agent attempts to construct an executable counterexample that
confirms or refutes the static hypothesis.

The goal is not to make the LLM itself the checker. Instead, \tool{} turns the
model's language-and-code understanding into an evidence-producing workflow:
field-centric rules define what must be checked, syntax-aware retrieval
collects the relevant implementation context, consistency reasoning records
supporting and missing evidence, and validation tests turn high-risk static
judgments into observable behavior when possible. The result is a framework
that can process natural-language standards while still producing reproducible
conformance claims.

We evaluate \tool on real-world implementations covering TLS, MQTT, and QUIC.
\tool generated 204 conformance-deviation reports.
Maintainers have confirmed and fixed 126 of them, roughly nine reports were
false positives, and 69 are still awaiting response or final adjudication. Two
confirmed reports correspond to previously unknown CVEs. The analysis is also
inexpensive relative to manual audit: using GPT-5.4 pricing as the cost model,
\tool spends \$1.122 per submitted report on average. These results suggest
that LLMs can support protocol conformance checking when their judgments are
grounded in field-level rules, implementation context, and validation
artifacts.

This paper makes the following contributions:
\begin{itemize}[leftmargin=*]
  \item We formulate protocol conformance checking as a field-centric
  specification-to-implementation consistency problem, using \ruleform{} rules
  to represent checkable obligations from natural-language standards.

  \item We design \tool, an evidence-driven pipeline that extracts field-level
  rules, aligns them to implementation variables and code contexts, reasons
  over consistency with explicit missing-evidence states, and validates
  suspected inconsistencies with executable tests when possible.

  \item We provide an empirical study, ablations, and case studies over
  real-world protocol implementations, producing 204 submitted reports, 126
  maintainer-confirmed fixes, roughly nine false positives, and two previously
  unknown CVEs.
\end{itemize}
